\documentclass[11pt]{article}
\usepackage[T1]{fontenc}
\usepackage{lmodern,microtype}
\usepackage{amsmath,amssymb,amsthm,mathtools}
\usepackage[letterpaper,margin=1in]{geometry}
\usepackage{booktabs,array,enumitem,longtable}
\usepackage{xcolor}
\usepackage[colorlinks=true,linkcolor=blue!55!black,citecolor=blue!55!black,urlcolor=blue!55!black]{hyperref}
\newcommand{\F}{F}
\newcommand{\K}{K}
\newcommand{\A}{\mathbb A}
\newcommand{\C}{\mathbb C}
\newcommand{\Z}{\mathbb Z}
\newcommand{\GL}{\operatorname{GL}}
\newcommand{\Irr}{\operatorname{Irr}}
\newcommand{\Nrd}{\operatorname{Nrd}}
\newcommand{\Nm}{\operatorname{N}}
\newcommand{\vol}{\operatorname{vol}}
\newcommand{\Gal}{\operatorname{Gal}}
\newcommand{\Aut}{\operatorname{Aut}}
\newcommand{\reg}{\mathrm{reg}}
\newcommand{\tH}{\widetilde H}
\newcommand{\tZ}{\widetilde Z}
\newcommand{\tT}{\widetilde T}
\newcommand{\Gb}{\overline G}
\newcommand{\GDb}{\overline{G_D}}
\newcommand{\Ptag}{\textnormal{\textbf{[P]}}}
\newcommand{\Etag}{\textnormal{\textbf{[E]}}}
\newcommand{\Ctag}{\textnormal{\textbf{[C]}}}
\newcommand{\Otag}{\textnormal{\textbf{[O]}}}
\theoremstyle{plain}
\newtheorem{theorem}{Theorem}[section]
\newtheorem{proposition}[theorem]{Proposition}
\newtheorem{lemma}[theorem]{Lemma}
\newtheorem{corollary}[theorem]{Corollary}
\theoremstyle{definition}
\newtheorem{definition}[theorem]{Definition}
\newtheorem{remark}[theorem]{Remark}
\newtheorem{example}[theorem]{Example}
\title{A metaplectic Jacquet--Langlands transfer for division algebras of prime degree, IV:\\
Local power maps, toral positivity, and conditional global comparison\\[0.4em]\normalsize Concluding Draft 6}
\author{David Joyner \and ChatGPT (OpenAI) \and Claude (Anthropic)}
\date{November 8, 2026}
\emergencystretch=3em
\begin{document}
\maketitle
\begin{abstract}
We conclude the fourth part of the metaplectic Jacquet--Langlands
series with a precise separation of proved local and arithmetic
results from two still-conditional goals.  We audit the ambient
cocycle and scalar-central parameter, recall the all-torus
good-element criterion, compute kernels and Haar moduli of
$n$th-power maps on finite etale tori modulo scalars, and
classify globally good regular rational classes.  Building on
Part III's unitary division-side character transform, we prove
that wild irreducible-character integrality is equivalent to
positive definiteness of every transferred irreducible character.
A new toral permutation-character theorem establishes positive
integral Fourier coefficients of the discriminant ratio for
\emph{every quadratic field extension} over a nonarchimedean
field of odd residue characteristic, for all odd $n$ including
$p\mid n$.  The complete numerical $r=2$, $n=3$ example and the
nonabelian gluing problem are developed separately in a
standalone companion paper, here begun as Draft~0.  We compare
our normalizations with Zou's correction of Flicker--Kazhdan,
state explicit conditions for split-place matching, and prove
a conditional geometric and trace-distribution comparison with
normalization factor $\Lambda$.  No full wild-place irreducible
correspondence, construction of exceptional local matching
functions, or unconditional ordinary-to-covered global
comparison is claimed.
\end{abstract}
\tableofcontents

\section{Introduction and mathematical status}\label{sec:intro}
The problem concerns three different local representation theories.
Write $H_v=D_v^\times$, and let $\GDb_v$ denote its specified $n$-fold
central cover; the cover $\Gb_v$ of $\GL_r(K_v)$ is defined inside the
same ambient cover.  An additional \emph{ordinary} group $H_v$ enters
through the transform $L_{D,v}$ of Part III.  These are not interchangeable.
In particular, a comparison between $\GDb$ and $\Gb$ has already been
obtained in Parts~I--II (here cited as Draft~8 of each) under their
stated test-function and spectral hypotheses.  We cite Part~III
as Draft~9 throughout.  The new proposed comparison between $H$ and $\GDb$ is a
different assertion.  This concluding version records the separate
rank-two wild case in a companion manuscript \cite{RankTwo0}.

\paragraph{Global assumptions.}
Let $K$ be a totally imaginary number field with $\mu_n\subset K$,
$n\geq1$ odd.  Let $r\geq2$ be prime, $(r,n)=1$, let $L/K$ be cyclic
of degree $r$, and let $D$ be a central division $K$-algebra of
degree $r$ containing $L$.  Put $S=\{v:D_v\text{ is division}\}$.
Fix integers $r',c$ such that
\begin{equation}\label{eq:parameters}
 rr'\equiv1\pmod n,\qquad r-1+2c\equiv0\pmod n,
 \qquad c_v\equiv r'c\pmod n.
\end{equation}
Thus $r-1+2rc_v\equiv0\pmod n$.  In fact $c_v$ is independent
of $v$ with these choices; the subscript recalls Part I's notation.
Locally $F=K_v$ is a nonarchimedean field containing $\mu_n$;
$p$ denotes its residue characteristic.  We call $v$ \emph{wild}
if $p\mid n$.

\paragraph{Status convention.}
\textbf{[P]} denotes a proof included here (possibly invoking
\textbf{[E]} inputs established in cited publications or earlier parts);
\textbf{[O]} denotes a proof outlined but not fully established, and
\textbf{[C]} denotes a conjectural goal requiring new mathematical
input.  When a proved proposition has explicit unresolved hypotheses,
we mark its formal implication [P] and mark the outstanding hypotheses
[C] separately, in agreement with the status usage of Part III,
Draft~9.

\begin{center}\small
\begin{tabular}{@{}p{0.49\linewidth}p{0.37\linewidth}p{0.07\linewidth}@{}}
\toprule
Assertion & Precise basis & Status\\\midrule
Diagonal and nonsplit good-element criterion, including wild $p$ & [FK, \S3]; Part I, \S3 & [E]\\
Compatibility of the Part I cocycle with $\sigma_{c_v}$ & Part I, Prop. 2.3(b); direct calculation & [P]\\
Power-map kernel, index, and Haar modulus on $T_E(F)/F^\times$ & Section \ref{sec:localpowers} & [P]\\
Global good-regular class criterion & Section \ref{sec:globalgood}; Part III Hasse principle & [P]\\
Division-side unitary character transform and tame integral case & Part III, Thms. 4.3, 7.6 & [E]\\
Quadratic toral permutation-character theorem (all odd $n$) & Theorem \ref{thm:quad-toral}; detailed example in companion paper & [P]\\
Wild irreducible-character integrality for $\widetilde{D^\times}$ & Section \ref{sec:wild} & [C]\\
Covered $D$--covered $\GL$ global simple trace comparison & Parts I--II under their hypotheses & [E]\\
Covered $D$--ordinary $D$ global comparison & Conditional implication in \S\ref{sec:geometric}; full matching open & [P]/[C]\\
Removal of auxiliary support and packet hypotheses & Section \ref{sec:spectral} & [C]\\
\bottomrule
\end{tabular}
\end{center}

\subsection*{The running example}
The series' running example is $K=\mathbb{Q}(\zeta_3)$, $n=3$, $r=2$, $D=(-1,4-\zeta_3)_K$, which is ramified exactly at $S=\{\mathfrak p_3,\mathfrak p_7\}$, with $\mathfrak p_3=(1-\zeta_3)$ and $\mathfrak p_7=(3+\zeta_3)$ (Part~I, \S1, and \texttt{notes/mjl\_running\_example.tex} in the project repository). Its wild place $\mathfrak p_3$ has $K_v=\mathbb{Q}_3(\zeta_3)$, and there $D$ is the quaternion division algebra treated in Remark~\ref{prop:all-three-tori} and in the companion paper \cite{RankTwo0}. Theorem~\ref{thm:quad-toral} gives the decompositions on all three quadratic tori at $\mathfrak p_3$. Theorem~\ref{thm:globalgood} describes the globally good classes, and Example~\ref{ex:ranktwoledger} is the global normalization ledger for this $D$, in which $\Lambda=1$ forces $\lambda_\infty=1/3$. The smallest open case of wild integrality is whether $L_D\mathbf 1$ at $\mathfrak p_3$ is irreducible of dimension $27$.

\section{The exact cocycle convention and star map}\label{sec:cocycle}
This section fixes the one place where a parameter mismatch could
invalidate every subsequent constant.  The ambient construction is
a restriction of a cocycle over $L_v$, whereas [FK] uses the standard
$\GL_r(F)$ Kazhdan--Patterson convention.

\subsection{Restriction and relabeling}
Let $\sigma_F$ be the parameter-zero Kazhdan--Patterson cocycle with
restriction to the diagonal torus
\begin{equation}\label{eq:parameterzero}
\sigma_F(\operatorname{diag}(a_i),\operatorname{diag}(b_i))
=\prod_{i<j}(a_i,b_j)_{n,F}.
\end{equation}
The $\GL$ subgroup of the ambient cover in Part I has, \emph{as an
extension}, cocycle model
\begin{equation}\label{eq:before}
  \sigma_F(g,h)^r\,(\det g,\det h)_{n,F}^{c}
  \quad\text{before relabeling }\mu_n.
\end{equation}
This assertion is Part I, Proposition 2.3(b).  Applying the
kernel relabeling $\zeta\mapsto\zeta^{r'}$ yields the model
\begin{equation}\label{eq:after}
 \sigma_{c_v}(g,h)=\sigma_F(g,h)
 (\det g,\det h)_{n,F}^{c_v}.
\end{equation}
The two parameters $c$ and $c_v$ must be kept separate: $c$
occurs before relabeling and $c_v=r'c\pmod n$ afterward.
The isomorphism in \eqref{eq:before} is an isomorphism of extensions,
not a claim that arbitrary choices of cocycle sections give
literally equal functions on all of $\GL_r(F)^2$.

\begin{proposition}[\Ptag\ Diagonal commutator audit]\label{prop:diag}
For $t=\operatorname{diag}(t_1,\ldots,t_r)$ and
$u=\operatorname{diag}(u_1,\ldots,u_r)$ in $\GL_r(F)$,
\begin{equation}\label{eq:matrix}
 [\widetilde t,\widetilde u]
 =\prod_{i,j}(t_i,u_j)_{n,F}^{A_{ij}},\qquad
 A=-I_r+(1+2c_v)J_r.
\end{equation}
Here $J_r$ is the all-ones matrix.  Its row sums are
$r-1+2rc_v\equiv0\pmod n$, so scalar lifts commute with the lifted
split torus.  The radical on $(F^\times/F^{\times n})^r$ consists
precisely of diagonal classes.
\end{proposition}
\begin{proof}
The commutator is the ratio
$\sigma_{c_v}(t,u)/\sigma_{c_v}(u,t)$.
The Hilbert symbol satisfies $(a,b)_n=(b,a)_n^{-1}$; this is
\emph{antisymmetry}, valid for all $n$.  The determinant twist
contributes $2c_v$ to each entry of $A$, and the parameter-zero
factor contributes $1$ off the diagonal.  This gives \eqref{eq:matrix}.
The Hilbert pairing on $X=F^\times/F^{\times n}$ is perfect.
As $X$ is killed by $n$ and $(r,n)=1$, the decomposition
$X^r=\Delta X\oplus\{(x_i):\sum x_i=0\}$ holds by averaging with
$r^{-1}$ modulo $n$.  On these summands $A$ acts, respectively, by
$r-1+2rc_v=0$ and $-1$.  Therefore $\ker A=\Delta X$.
\end{proof}

\begin{remark}
The assumption $(r,n)=1$ also follows algebraically from
$n\mid r-1+2rc_v$, because any common divisor of $r$ and $n$ must
divide $1$.  We retain it because it is a global hypothesis of the
series.  The oddness of $n$ is used for the ordinary star-map
convention below, not for Hilbert-symbol antisymmetry.
\end{remark}

\subsection{Star and central characters}\label{subsec:star}
For an element $h$ of the ordinary group, choose any lift
$\widetilde h$ and define $h^*=\widetilde h^{\,n}$.  It is
independent of the lift because $\mu_n$ is central.  In the present
scalar-central cover, for $a\in F^\times$,
\begin{equation}\label{eq:starcentral}
 (ah)^*=a^*h^*,\qquad a^*=\widetilde a^{\,n}.
\end{equation}
If $h_1,h_2$ commute in an ordinary torus, their lifted commutator
is in $\mu_n$, and the class-two power identity
\begin{equation}\label{eq:starpower}
 (\widetilde h_1\widetilde h_2)^n
 =\widetilde h_1^{\,n}\widetilde h_2^{\,n}
 [\widetilde h_2,\widetilde h_1]^{n(n-1)/2}
 =\widetilde h_1^{\,n}\widetilde h_2^{\,n}
\end{equation}
uses $n$ odd.  Thus star is multiplicative on the lifted torus.
This agrees with the convention in [FK, \S4] and [Zou, \S3]:
Zou writes $x^*=u(x)s(x)^n$, where $u(x)=\pm1$ is the
correction defined in \cite[\S4]{FK} and \cite[\S3.1]{Zou24}.
For odd $n$ the extra sign is absent in the
Flicker--Kazhdan convention: the star is the $n$th power of a lift.
The correction $u(x)$ is introduced for the even-$n$ convention;
we do not infer that an arbitrary sign automatically lies in
$\mu_n$.  In the odd-degree setting, therefore, $u(x)=1$
by definition of the relevant convention (compare
\cite[\S4]{FK} and \cite[\S3.1]{Zou24}).  For a scalar $\zeta\in\mu_n(F)$,
$\zeta^*=1$ in the standard diagonal cocycle model: the
self-symbol $(\zeta,\zeta)_n=1$ for odd $n$, so
$s(\zeta I_r)^n=s(I_r)=1$.  This last statement is also
invariant under the extension isomorphisms in Part I.
For representation characters we first fix a genuine unitary
character $\widetilde\theta$ of the central preimage $\tZ$ and
\emph{then} put $\theta(a)=\widetilde\theta(a^*)$.
The reverse implication is not unique: $\theta$ alone does not
specify $\widetilde\theta$.  For $\zeta\in\mu_n(F)$ we have
$\zeta^*=1$, hence $\theta(\zeta)=1$.  Globally the rational
splitting gives $a^*=\iota(a^n)$ for $a\in K^\times$;
if $\widetilde\theta$ is automorphic, it is trivial on the rational
splitting, so $\theta(a)=1$ on $K^\times$.
Test functions use the inverse central characters
$\theta^{-1}$ and $\widetilde\theta^{-1}$.

\section{Good elements for every maximal torus}\label{sec:good}
A regular maximal torus in $\GL_r(F)$ may be nonsplit even at a
place where $D_v\simeq M_r(F)$.  This is the situation not treated by
the diagonal matrix alone.  The following calculation is not new:
it is the proof of [FK, \S3, Proposition] in the notation of this
series, and also underlies Part I, Proposition 3.5.

\begin{theorem}[\Etag\ Flicker--Kazhdan good-element theorem]\label{thm:FKgood}
Let $E=\prod_{i=1}^kE_i$ be a finite \emph{\'etale} $F$-algebra
of dimension $r$, embedded as a regular maximal torus
$T_E(F)=E^\times$ in $\GL_r(F)$.  Under \eqref{eq:parameters},
for any $t\in T_E(F)$ and any lift $\widetilde t$,
\begin{equation}\label{eq:FKtorus}
 [\widetilde t,\widetilde u]
 =(\Nm_{E/F}t,\Nm_{E/F}u)_{n,F}^{1+2c_v}
   \prod_{i=1}^k(t_i,u_i)_{n,E_i}^{-1}.
\end{equation}
The radical of this commutator pairing is
\begin{equation}\label{eq:everygood}
 \{t\in T_E(F):[\widetilde t,\widetilde u]=1\ \forall u\in T_E(F)\}
 =F^\times T_E(F)^n.
\end{equation}
For regular $t$, this is exactly the criterion that $\widetilde t$
is good.  Both identities are valid when $p\mid n$.
\end{theorem}
\begin{proof}[Proof from the cited formula, included for clarity]
The commutator formula is [FK, \S3], using [KP, Prop.~0.1.5];
its restriction to a single field is exactly Part I,
Proposition 3.1 after relabeling \cite{JoynerI}.  Write $N=\Nm_{E/F}$.
Norm compatibility of local Hilbert symbols gives
\[
 (Nt,Nu)_{n,F}
 =\prod_i(Nt,u_i)_{n,E_i}.
\]
Thus \eqref{eq:FKtorus} equals
\[
 \prod_i\bigl((Nt)^{1+2c_v}t_i^{-1},u_i\bigr)_{n,E_i}.
\]
The nondegeneracy of the $n$th Hilbert symbol over each $E_i$
(including wildly ramified $E_i$) implies that the commutator is
trivial for all $u$ exactly when
\begin{equation}\label{eq:pointwisegood}
 t_i\equiv(Nt)^{1+2c_v}\pmod{E_i^{\times n}}
 \quad (1\leq i\leq k).
\end{equation}
This directly says $t\in F^\times E^{\times n}$.  Conversely,
if $t=ay^n$ for $a\in F^\times$, then
$Nt=a^r(Ny)^n$; therefore
$(Nt)^{1+2c_v}\equiv a^{r(1+2c_v)}\equiv a\equiv t_i$
modulo $E_i^{\times n}$ by scalar centrality.  This proves the
criterion.  No tameness assumption has been introduced.
\end{proof}

\begin{example}[The diagonal rank-two check]\label{ex:diag}
Take $r=2$, $n=3$ and $F=\mathbb Q_3(\zeta_3)$.
In the notation of Part I, Example 3.7, choose $c=1$, so
$c_v\equiv2\pmod3$.  Formula \eqref{eq:matrix} becomes
\[
 A\equiv\begin{pmatrix}1&-1\\-1&1\end{pmatrix}\pmod3,
 \qquad
 [\widetilde t,\widetilde u]
 =(t_1/t_2,u_1/u_2)_{3,F}.
\]
Thus $\operatorname{diag}(t_1,t_2)$, when regular, is good exactly
when $t_1/t_2\in F^{\times3}$.  If a quadratic field $E/F$ is
embedded as a nonsplit maximal torus, the same theorem says that
goodness is equivalent to $t\in F^\times E^{\times3}$; the two
tests are instances of the same formula, not separate conjectures.
\end{example}

\begin{remark}[Why goodness is not enough for a regular-root formula]\label{rem:power-singular}
The power of a regular element in a \emph{split} torus need not be
regular.  For $n=3$ and $F=\mathbb Q_3(\zeta_3)$, the element
$\operatorname{diag}(1,\zeta_3)$ is regular but not good:
$\zeta_3\notin F^{\times3}$, since $F$ does not contain $\zeta_9$.
Its cube is the identity.  A character or orbital-integral formula with
$\Delta(t^n)$ in the denominator must restrict to $t^n$ regular
(as in Zou's equation (3.2)), or use a separately proved germ extension.
On a degree-$r$ \emph{field} torus, by contrast, if $t$ is regular
and $t^n\in F$, then Kummer theory gives
$[F(t):F]\mid n$; since $[F(t):F]=r$ and $(r,n)=1$, this is
impossible.  Because $r$ is prime, $F(t^n)=F(t)$ and $t^n$ is regular.
\end{remark}

\section{Local power-map calculation for \texorpdfstring{\'etale}{etale} tori}\label{sec:localpowers}
The key local distinction is between a field centralizer (compact
modulo scalars) and a product of fields (not compact modulo scalars).
Both occur in the split \emph{group}.  The next proposition gives
an exact count of the roots and the size of the image at wild and
tame places, including the factors that a global trace identity
must account for.

\begin{lemma}[\Ptag\ Kummer injection for a product algebra]\label{lem:Kumprod}
Let $E=\prod_{i=1}^kE_i$ be finite \emph{\'etale} of dimension $r$
over $F$, with $\mu_n\subset F$ and $(r,n)=1$.  Then
\begin{equation}\label{eq:Kummerinj}
 F^\times/F^{\times n}\longrightarrow
 E^\times/E^{\times n}
\end{equation}
is injective.
\end{lemma}
\begin{proof}
Suppose $a\in F^\times$ is an $n$th power in every $E_i$.
The Kummer extension $F(a^{1/n})/F$, since $\mu_n\subset F$,
is cyclic of degree $d\mid n$ and embeds in every $E_i$.
Therefore $d\mid[E_i:F]$ for every $i$, so also
$d\mid\sum_i[E_i:F]=r$.  Coprimality gives $d=1$ and
$a\in F^{\times n}$.
\end{proof}

\begin{proposition}[\Ptag\ Kernel, index, and Haar modulus]\label{prop:power-etale}
Under the hypotheses of Lemma \ref{lem:Kumprod}, put
$A_E=E^\times/F^\times$ and $P_n([t])=[t^n]$.  Then
\begin{align}
 \ker(P_n)&\simeq \mu_n(E)/\mu_n(F)
      \simeq \mu_n^{\,k}/\mu_n,
 &|\ker(P_n)|&=n^{k-1},\label{eq:kernel}\\
 [A_E:P_n(A_E)]&=n^{2(k-1)}|n|_F^{-(r-1)}.
 \label{eq:index}
\end{align}
In particular, $P_n(A_E)$ is open and closed.  With \emph{any}
Haar measure $da$ on $A_E$ and its restriction to $P_n(A_E)$,
for compactly supported integrable $q$ on that image,
\begin{equation}\label{eq:haarJac}
 \int_{A_E}q(a^n)\,da
 =n^{k-1}|n|_F^{-(r-1)}
   \int_{P_n(A_E)}q(b)\,db.
\end{equation}
\end{proposition}
\begin{proof}
If $t^n=a\in F^\times$, then $a$ is an $n$th power in every
$E_i$; by Lemma \ref{lem:Kumprod}, write $a=b^n$ with $b\in F^\times$.
Thus each $t_i/b\in\mu_n(E_i)=\mu_n(F)$, giving
\eqref{eq:kernel}, including its canonical description.

For a local field $M$ containing $\mu_n$ the Kummer quotient has
\begin{equation}\label{eq:localKumIndex}
 |M^\times/M^{\times n}|=n^2|n|_M^{-1}.
\end{equation}
Indeed the discrete valuation contributes $n$; the finite
torsion unit factor containing $\mu_n$ contributes $n$; and
the torsion-free $p$-adic unit factor contributes
$|n|_M^{-1}$ (multiplication by $n$ on its $\Z_p$-lattice).  Alternatively, \eqref{eq:localKumIndex}
follows from local duality.
Since $F^\times/F^{\times n}$ injects into
$E^\times/E^{\times n}$, the group $A_E/A_E^n$ is its cokernel:
\[
 |A_E/A_E^n|
 =\frac{\prod_i n^2|n|_{E_i}^{-1}}
             {n^2|n|_F^{-1}}
 =n^{2(k-1)}|n|_F^{-(r-1)},
\]
where $|n|_{E_i}=|n|_F^{[E_i:F]}$ for normalized absolute values
on local fields.  This proves \eqref{eq:index}.

For the Haar statement, write $K_0=\ker(P_n)$, of order
$n^{k-1}$.  Choose a sufficiently small compact open subgroup
$U\subset A_E$ consisting of principal-unit classes and having
$U\cap K_0=\{1\}$.  The $p$-adic logarithm identifies $U$ with
a free $\mathbb Z_p$-module of rank
$(r-1)[F:\mathbb Q_p]$: the factors $E_i^\times$ have combined
principal-unit rank $r[F:\mathbb Q_p]$, and quotienting by the
scalar principal units subtracts $[F:\mathbb Q_p]$.
Multiplication by $n$ maps $U$ injectively into itself, with
\[
 [U:U^n]=p^{v_p(n)(r-1)[F:\mathbb Q_p]}
        =|n|_F^{-(r-1)}.
\]
Since $P_n^{-1}(U^n)=K_0U$ and $K_0\cap U=\{1\}$, its
pushforward of Haar measure gives
\[
 \frac{\operatorname{vol}(P_n^{-1}(U^n))}
      {\operatorname{vol}(U^n)}
   =\frac{|K_0|\operatorname{vol}(U)}
          {\operatorname{vol}(U)/[U:U^n]}
   =n^{k-1}|n|_F^{-(r-1)}.
\]
The pushforward is translation-invariant on the open image
$P_n(A_E)$, so this ratio is its constant density on that image,
proving \eqref{eq:haarJac}.  Notice that this argument requires
neither a global quotient-index formula nor a choice of scale for
Haar measure.
\end{proof}

\begin{remark}[A measure-theoretic warning]\label{rem:haar}
For a general open endomorphism $\varphi:A\to A$ of a locally
compact abelian group with finite kernel, its pushforward Haar density
need not equal $[A:\varphi A]/|\ker\varphi|$.  For example, on
$A=\mathbb Z$ with counting measure, $a\mapsto na$ has density
$1$ on $n\mathbb Z$, not $n$.  The correct local computation here is
$|\ker P_n|[U:U^n]$, where $U$ is the compact open subgroup in
the proof.  For this specific $A_E$ it happens to simplify to
$[A_E:P_nA_E]/|\ker P_n|$; that numerical equality is not a
general change-of-variables theorem.
\end{remark}

\begin{example}[Three local torus types]
If $E/F$ is a field of degree $r$, then $k=1$ and $P_n$ is
injective, with index and Haar modulus $|n|_F^{-(r-1)}$.
If $E=F^r$ then $k=r$ and the kernel has order $n^{r-1}$.
For $r=2,n=3$, $F=\mathbb Q_3(\zeta_3)$ and $|3|_F=1/9$,
\[
\begin{array}{c|ccc}
\text{type}&|\ker P_3|&[A_E:A_E^3]&\text{factor in \eqref{eq:haarJac}}\\\hline
E/F\text{ quadratic}&1&9&9\\
E=F\times F&3&81&27
\end{array}
\]
These integers quantify why one cannot use the division-torus
substitution formula unchanged on the full split group.
\end{example}

\section{A global classification of good regular classes}\label{sec:globalgood}
The local theorem now applies to the actual torus at \emph{every}
place.  Combining it with the Hasse principle of Part III removes
the local-to-global ambiguity for power classes of a degree-$r$
field, though not the independent measure and spectral questions.

\begin{proposition}[\Etag\ Almost-everywhere powers modulo scalars]\label{prop:Hasse}
Let $E/K$ be a field extension of degree $r$ with $(r,n)=1$ and
$\mu_n\subset K$.  For every finite set $\Sigma$ of places,
\begin{equation}\label{eq:Hasse}
 E^\times/(K^\times E^{\times n})\longrightarrow
 \prod_{v\notin\Sigma}
 \frac{(E\otimes_K K_v)^\times}
 {K_v^\times(E\otimes_KK_v)^{\times n}}
\end{equation}
is injective.  This is Part III, Proposition 11.3 \cite{JoynerIII}.
\end{proposition}
\begin{proof}[Recalled proof]
Let $t$ be locally scalar times an $n$th power outside $\Sigma$.
Choose $u,k\in\Z$ with $ru+nk=1$ and put
$a=N_{E/K}(t)^u\in K^\times$.  If $t=a_v y_v^n$ in the local
\'etale algebra, taking norms gives the identity
\[
 t/a=\bigl(a_v^k y_v/N_{E_v/K_v}(y_v)^u\bigr)^n.
\]
Thus $t/a$ is an $n$th power in every completion of $E$ over all
$v\notin\Sigma$.  If its class in $E^\times/E^{\times n}$ were
nontrivial, the nontrivial Kummer extension
$E((t/a)^{1/n})/E$ would, by Chebotarev, admit infinitely many
unramified places at which $t/a$ is not an $n$th power.
This contradicts the hypothesis.  Hence $t\in K^\times E^{\times n}$.
\end{proof}

\begin{theorem}[\Ptag\ Global good-regular criterion]\label{thm:globalgood}
Assume the covers and rational splittings of Part I.  Let
$G$ denote $D^\times$ or $\GL_r$, and let $\gamma\in G(K)$ be
regular elliptic, so $E=K[\gamma]$ is a field of degree $r$.
The following are equivalent:
\begin{enumerate}[label=(\alph*)]
\item the rational lift $\iota(\gamma)$ is good at every place;
\item it is good at every place outside some finite set $\Sigma$;
\item $\gamma\in K^\times E^{\times n}$.
\end{enumerate}
The map $[s]\mapsto[s^n]$ is a bijection from the noncentral
classes in $E^\times/K^\times$ onto the globally good regular
classes belonging to $E$.  Furthermore
\begin{equation}\label{eq:Wstabilizer}
 W_{s^n}=W_s,\qquad
 W_s=\{\sigma\in\Aut(E/K):\sigma(s)\in K^\times s\}.
\end{equation}
\end{theorem}
\begin{proof}
At each finite $v$, the ordinary centralizer is
$T_v=(E\otimes_K K_v)^\times$.  If $D_v$ splits, apply
Theorem \ref{thm:FKgood}, transported through Part I, Proposition~2.6 (the split-place
isomorphism of the specified covers).  If $D_v$ is division, use Part I,
Proposition 3.5(a), whose commutator is the same by Part I,
Proposition 3.1.  At complex places the $n$th-power map on
$T_v$ is surjective.  Thus at every place the local goodness
condition is precisely
$\gamma\in K_v^\times T_v^n$.
Implications (c)$\Rightarrow$(a)$\Rightarrow$(b) follow,
while (b)$\Rightarrow$(c) is Proposition \ref{prop:Hasse}.

The power map $E^\times/K^\times\to E^\times/K^\times$
is injective: if $s^n\in K^\times$, Kummer theory shows
$[K(s):K]\mid n$, while $[K(s):K]\mid [E:K]=r$.
Hence $s\in K^\times$.  Because $r$ is prime, every noncentral
$s\in E$ generates $E$, and the same is true of $s^n$.
This proves the bijection.  Finally
$\sigma(s^n)/s^n\in K^\times$ if and only if
$(\sigma(s)/s)^n\in K^\times$, which is equivalent by injectivity
to $\sigma(s)/s\in K^\times$.  Equation \eqref{eq:Wstabilizer}
follows.
\end{proof}

\begin{remark}[Precise consequence for the geometric side]
Part I, Proposition 6.1, uses the rational regular classes which
are good everywhere, with coefficient
\begin{equation}\label{eq:globalvol}
 c(\gamma)=\frac{\vol\bigl(\A_K^\times E^\times\backslash
 \A_E^\times;dt\bigr)}{|W_\gamma|}.
\end{equation}
Theorem \ref{thm:globalgood} parametrizes its noncentral
summation set by global $[s^n]$ and preserves $W_\gamma$.
It does \emph{not} replace the local root sums in Zou's
character formula: the global root class may be unique even when
a localized split torus has $n^{k_v-1}$ roots modulo scalars.
\end{remark}

\section{Wild local integrality: exact reductions and a toral theorem}\label{sec:wild}
The division-side transform of Part~III is an isometric isomorphism of
spaces of invariant class functions.  Wild-place irreducible-character
integrality does not follow from this analytic statement.  We record
here the exact character-theoretic criterion and a positive result on
quadratic field tori.  Detailed ramified and unramified computations
for $F=\mathbb Q_3(\zeta_3)$ belong to the separate companion project
\emph{Wild-place metaplectic character transfer in degree two},
Draft~0, cited below as \cite{RankTwo0}.

\subsection{Normalization and the unitary character matrix}\label{subsec:integral}
Let $F$ be a nonarchimedean field containing $\mu_n$, let $D/F$
be a division algebra of prime degree $r$ with $(r,n)=1$, and put
$H=D^\times$ and $Z=F^\times$.  Write $\tH$ for the scalar-central
cover specified in Part~I.  For a regular element $t$ in a maximal
field $E\subset D$, put
\begin{equation}\label{eq:Delta}
\Delta(t)=\left|\det(1-\operatorname{Ad}(t)\mid
\mathfrak h/\mathfrak t)\right|_F^{1/2},\qquad
R_E(t)=\frac{\Delta(t)}{\Delta(t^n)}.
\end{equation}
The character transform of Part~III, with compatible central
characters $\theta,\widetilde\theta$, is
\begin{equation}\label{eq:LD}
(L_D\Phi)(z t^*)=b_D\widetilde\theta(z)R_E(t)\Phi(t),
\qquad b_D=|n|_F^{-(r-1)/2};
\end{equation}
it vanishes on nongood regular lifts.  Its continuously extended
identity value is
\begin{equation}\label{eq:cD}
(L_D\Theta_\rho)(1)=c_D\dim\rho,\qquad
c_D=|n|_F^{-(r^2-1)/2}.
\end{equation}
Theorem~4.3 and Corollary~5.1 of Part~III prove that $L_D$ is a
\emph{unitary onto map} between the fixed-central-character invariant
$L^2$ spaces and admit finite expansions
\begin{equation}\label{eq:matrixchars}
L_D\Theta_\rho=\sum_\tau a_{\rho\tau}\Theta_\tau,
\qquad a_{\rho\tau}=\langle L_D\Theta_\rho,\Theta_\tau\rangle.
\end{equation}
We do not assume that $a_{\rho\tau}$ is integral at wild places.

\begin{proposition}[\Ptag\ Integrality--positivity criterion]\label{prop:int}
For the entire unitary matrix $A=(a_{\rho\tau})$, the following
conditions are equivalent:
\begin{enumerate}[label=(\roman*)]
\item $L_D$ bijects the irreducible-character bases with positive signs;
\item $a_{\rho\tau}\in\mathbb Z$ for all $\rho,\tau$;
\item $a_{\rho\tau}\ge0$ is real for all $\rho,\tau$;
\item every $L_D\Theta_\rho$ is a continuous positive-definite
function on $\tH$.
\end{enumerate}
Positive definiteness for \emph{one} row is not sufficient for
integrality of that row.
\end{proposition}
\begin{proof}
A row with integral entries and squared norm one has exactly one
nonzero entry, equal to $\pm1$.  Equation~\eqref{eq:cD} rules out the
minus sign, and unitarity and surjectivity give the bijection.
Conversely, permutation of the character bases gives integrality.
For nonnegative entries, orthogonality of distinct rows implies
that their supports are disjoint.  Every column has squared norm
one; it meets at most one row, and its sole nonzero entry therefore
equals one.  Each row has norm one, so it also contains exactly one
entry equal to one.  Finally, Peter--Weyl orthogonality for the
compact-mod-center group with fixed unitary central character
identifies positive-definite invariant functions with those having
nonnegative coefficients in the irreducible-character basis.
A single nonnegative norm-one row may nevertheless contain several
nonintegral entries (for example $1/\sqrt2,1/\sqrt2$).
\end{proof}

\subsection{The toral Fourier reduction}
For a maximal field $E/F$, $A_E=E^\times/F^\times$ is compact
abelian.  The continuous discriminant ratio is locally constant and
hence has a finite Fourier expansion
\begin{equation}\label{eq:R-Fourier}
R_E(t)=\sum_{\chi\in\widehat{A_E}}m_\chi\chi(t).
\end{equation}
A genuine central character on $\widetilde{E^\times}$ compatible
with an ordinary character $\alpha$ identifies the coefficient
$m_\chi$ with the multiplicity of the corresponding genuine
Stone--von Neumann torus representation in the \emph{restricted}
transfer $\mathcal L_E\alpha$ (Part~I, Proposition~3.5,
and Part~III).  This is a statement about the full preimage
of the field torus; it is not an assertion that the torus modules
extend to $\tH$.  Consequently,
\begin{equation}\label{eq:toralnecessary}
 m_\chi\in\mathbb Z_{\ge0}\quad\text{for all $\chi$}
\end{equation}
is sufficient for positive integral \emph{torus restrictions} of
all transferred ordinary characters.  It does not imply any
$a_{\rho\tau}\in\mathbb Z$ for the full nonabelian cover.

\begin{theorem}[\Ptag\ Quadratic toral permutation positivity]
\label{thm:quad-toral}
Let $F/\mathbb Q_p$ be finite with $p$ odd, let $n$ be odd with
$\mu_n\subset F$, and let $E/F$ be a quadratic field extension.
Then $R_E$ is a permutation character of a finite quotient of
$E^\times/F^\times$.  In particular the Fourier coefficients in
\eqref{eq:R-Fourier} belong to $\mathbb Z_{\ge0}$, even when $p\mid n$.
At the identity, $R_E(1)=|n|_F^{-1}$.
\end{theorem}
\begin{proof}
Write $n=p^a m$ with $(m,p)=1$ and set
$G=E^1=\ker N_{E/F}$.  Hilbert~90 identifies $E^\times/F^\times$
with $G$ by $t\mapsto u=t/\sigma(t)$.
If $a=0$ the ratio is identically one, so suppose $a>0$.
Set $G_k=\{u\in G:v_E(1-u)\ge k\}$ and $G_0=G$.
A root $\zeta\in\mu_n$ with nontrivial prime-to-$p$ component
cannot have $\zeta u\equiv1$ modulo the maximal ideal: reduction of
$u$ has norm one and its intersection with the scalars is $\{\pm1\}$.
For $\zeta\in\mu_{p^a}\setminus\{1\}$, the ultrametric inequality,
together with $\sigma(u)=u^{-1}$ and $p$ odd, gives
\[
v_E(1-\zeta u)=\min\bigl(v_E(1-u),v_E(1-\zeta)\bigr).
\]
Indeed, cancellation at equal valuations would give both
$u\equiv\zeta^{-1}$ and $u^{-1}\equiv\zeta^{-1}$ beyond the
valuation of $1-\zeta$, contradicting
$v_E(\zeta-\zeta^{-1})=v_E(1-\zeta)$.
Factoring $1-u^n$ yields
\begin{equation}\label{eq:toralradial}
R_E(u)=q_E^{\frac12\sum_{\zeta\in\mu_{p^a}\setminus\{1\}}
\min(v_E(1-u),v_E(1-\zeta))},
\end{equation}
where $q_E$ is the size of the residue field of $E$.
This formula also defines the continuous central value.

Put $t_\zeta=v_E(1-\zeta)$,
$T=\max_{\zeta\ne1}t_\zeta$ and
$M_k=q_E^{\frac12\sum_\zeta\min(k,t_\zeta)}$ ($M_0=1$).
The $M_k$ are integers, since $\zeta$ and $\zeta^{-1}$ occur
in pairs, and $M_k=M_T$ for $k\ge T$.
If $E/F$ is unramified and $q=q_F$, then $q_E=q^2$ and
\[
[G:G_k]=(q+1)q^{k-1}\quad(k\ge1).
\]
Writing $s_k=\#\{\zeta\ne1:t_\zeta\ge k\}$, we have
$M_k/M_{k-1}=q^{s_k}$ and $s_k$ is even.  The quantities
\[
c_k=\frac{M_k-M_{k-1}}{(q+1)q^{k-1}}
\quad (1\le k\le T)
\]
are nonnegative integers: $q+1\mid q^{s_k}-1$, and when
$s_k>0$, the preceding $M_{k-1}$ is divisible by $q^{k-1}$.
Therefore
\[
R_E=\mathbf1_G+\sum_{k=1}^{T}c_k[G:G_k]\mathbf1_{G_k}.
\]
Each summand $[G:G_k]\mathbf1_{G_k}$ is a permutation
character of $G/G_k$.

If $E/F$ is ramified, then $q_E=q$, the $t_\zeta$ are even,
and $G_{2j}=G_{2j+1}$, with
$[G:G_1]=2$ and $[G:G_{2j}]=2q^j$ for $j\ge1$.
Set
\[
c_1=\frac{M_1-1}{2},\qquad
c_{2j}=\frac{M_{2j+1}-M_{2j-1}}{2q^j}.
\]
The first is a nonnegative integer because $M_1$ is odd.
When the numerator of $c_{2j}$ is nonzero, some pair
$\zeta,\zeta^{-1}$ has $t_\zeta\ge2j$; hence
$q^j\mid M_{2j-1}$.  The quotient
$M_{2j+1}/M_{2j-1}$ is a power of odd $q$, so its difference
from one is even.  Thus every $c_{2j}\in\mathbb Z_{\ge0}$ and
\[
R_E=\mathbf1_G+c_1[G:G_1]\mathbf1_{G_1}
+\sum_{1\le j\le T/2}c_{2j}[G:G_{2j}]\mathbf1_{G_{2j}}.
\]
The norm-one filtration indices follow directly by writing
$E=F(\sqrt d)$, $v_F(d)=1$, and parameterizing
$G_1$ by $(1+c\sqrt d)/(1-c\sqrt d)$ with $c\in\mathcal O_F$;
congruence modulo $G_{2j}$ records $c$ modulo
$\mathfrak p_F^j$.  This completes the two cases.
At the center $R_E(1)=|n|_F^{-1}$, by the cyclotomic product.
\end{proof}

\begin{remark}[\Ptag\ A quantified rank-two instance]\label{prop:all-three-tori}
For $F=\mathbb Q_3(\zeta_3)$, $n=3$ and the quaternion
$D/F$, every quadratic maximal subfield $E/F$ satisfies
Theorem~\ref{thm:quad-toral}.  The detailed decompositions are
\[
R_{E_{\mathrm{ur}}}=1+2\operatorname{Reg}(C_4),\qquad
R_{E_{\mathrm{ram}}}
=1+\operatorname{Reg}(C_2)+\operatorname{Reg}(C_6).
\]
These give genuine torus restriction multiplicities
$(3,2,2,2)$ and $(3,2,1,1,1,1)$, respectively, of
degree-three constituents, hence total dimension $27$.
Both ramified quadratic isomorphism classes yield the same
formula.  The full derivations and the status of the unproved
nonabelian gluing step are recorded in \cite{RankTwo0}.
\end{remark}

\subsection{What the local reductions prove, and what they do not}
\label{subsec:strategywild}
Theorem~\ref{thm:quad-toral} establishes toral Fourier positivity
for \emph{all} quadratic extensions and odd covering degrees,
including wild primes, but does not establish positive
definiteness on the full cover.  The logical implications are
\[
\text{positive character on every covered torus}
\quad\not\!\!\Longrightarrow\quad
\text{positive-definite character on }\tH.
\]
At a fixed finite quotient this is visible even for
$C_2\times C_2$: the function equal to $2$ at the identity
and $-2$ at each of the three involutions restricts to twice
the sign character on every cyclic subgroup of order two, but
equals $-\chi_0+\chi_1+\chi_2+\chi_3$ on the full group, so it
is not positive definite.  (The function equal to $2$ at the
identity and $0$ elsewhere shows separately that positivity
does not force integrality: it restricts to the regular
character on each cyclic subgroup, but is half the sum of the
four irreducible characters.)
The finite-level Brauer/Clifford and induced-type strategies in
\cite{RankTwo0} address this additional \emph{nonabelian gluing}
requirement.  For arbitrary prime $r>2$, even the toral
Fourier-positivity step remains open.
No theorem here identifies a wild irreducible genuine image or
proves the new ordinary-$D^\times$ to covered-$D^\times$
automorphic comparison unconditionally.

\section{Connection with Zou's corrected character relation}\label{sec:Zou}
For the Kazhdan--Patterson cover of the \emph{split group},
Zou \cite[\S3]{Zou24} revisits the Flicker--Kazhdan metaplectic
correspondence and records a corrected scalar in the character
relation.  This is an essential input, but its objects are
$\widetilde{\GL_r(F)}$ and ordinary $\GL_r(F)$, not directly
$\tH$ and $H$ at a division place.

\begin{proposition}[\Etag\ Zou's root-sum constant in our case]\label{prop:Zouconstant}
Write $d_F(m)=m|m|_F^{-1/2}$ and
$d_r=\gcd(n,2rc_v+r-1)$.  Under \eqref{eq:parameters},
$d_r=n$.  In Zou's equation (3.2), the coefficient preceding
the root sum in rank $r$ is
\begin{equation}\label{eq:Zouprefactor}
 \frac{1}{d_F(d_r)|n|_F^{r/2}}
 =\frac{1}{n|n|_F^{(r-1)/2}}=\frac{b_D}{n}.
\end{equation}
For regular elliptic $x$ with degree-$r$ field centralizer,
$x^n$ is automatically regular by Remark~\ref{rem:power-singular}.
The equation $t^n=x^n$ then has exactly $n$ solutions
$t=\zeta x$, $\zeta\in\mu_n$.  If the linear character
$\pi$ has central character $\omega_\pi$ trivial on $\mu_n$,
these summands agree and the effective coefficient is $b_D$.
This restriction on $\omega_\pi$ is essential: without it, the
scalar sum is $\sum_{\zeta\in\mu_n}\omega_\pi(\zeta)=0$,
so equation (3.2) instead forces the covered character to vanish
at $x^*$.  For a genuine lift pair the compatibility
$\omega_\pi(a)=\omega_{\widetilde\pi}(a^*)$
and $\zeta^*=1$ imply $\omega_\pi|_{\mu_n}=1$ automatically.
\end{proposition}
\begin{proof}
The equation $d_r=n$ is immediate from scalar centrality;
\eqref{eq:Zouprefactor} follows by substitution.  Regularity of
$x^n$ follows from the field case of Remark~\ref{rem:power-singular}.
Any solution $t$ centralizes $x^n$, so $t\in F[x]^\times$.  In the field $F[x]$, $(t/x)^n=1$,
and $\mu_n(F[x])=\mu_n(F)$.  Thus $t=\zeta x$.
The Weyl discriminant is invariant under central scalar
multiplication, the ordinary character takes the same value
because it is trivial on $\mu_n$, and $(\zeta x)^*=x^*$
by the explicit odd-$n$ scalar star computation in
Section~\ref{subsec:star}.
These are precisely the equal terms in Zou's (3.2).
\end{proof}

\begin{remark}[What is and is not composed]
At the ramified places one may compare the \emph{desired}
$\GDb_v\leftrightarrow\Gb_v$ character relation of Parts I--II
with (a) the classical linear Jacquet--Langlands elliptic
character identity and (b) Zou's formula.  The common coefficient
$b_D$ and the familiar sign $(-1)^{r-1}$ are consistent with
this comparison.  Nevertheless \emph{the existence of an
irreducible genuine image on the covered division side is the
missing conclusion}; combining two known formulas cannot be
used to assume that missing representation exists.
For non-elliptic split centralizers, retain Zou's entire
root sum, including the central-character factors.
\end{remark}

\section{Global geometric comparison: proved input and new obligations}\label{sec:geometric}
There are two comparisons to keep distinct:
\begin{equation}\label{eq:twocomparisons}
 \begin{gathered}
 \GDb(\A_K)\ \longleftrightarrow\ \Gb(\A_K)
 \quad\text{(Parts I--II, under their hypotheses)},\\[3pt]
 D^\times(\A_K)\ \dashrightarrow\ \GDb(\A_K)
 \quad\text{(proposed additional comparison)}.
 \end{gathered}
\end{equation}
The first has a proved simple trace-formula comparison with
regular-support or later matching hypotheses (Part I, Theorem 6.6;
Part II, Theorem 6.2).
The second is not a formal consequence of the first or of the
unitarity of $L_D$.

\subsection{Geometric coefficients and the product formula}
For a regular elliptic rational $\gamma$, with $E=K[\gamma]$,
Part I identifies the torus coefficient $c(\gamma)$, recalled in
\eqref{eq:globalvol}. The same torus measure $dt$ must be used
whenever this term is compared with the corresponding linear
inner-form orbital integral.
Part I uses the normalized orbital integral
$J_v(\gamma,f)=\Delta_v(\gamma)O_v(\gamma,f)$.
For a rational regular element the product formula implies
\begin{equation}\label{eq:discprod}
 \prod_v\Delta_v(\gamma)=1.
\end{equation}
This is because each discriminant is the normalized absolute
value of a nonzero element of $K$ (or its square root), including
the determinant normalization in \eqref{eq:Delta}.
The local \emph{ratio} of discriminants
$\Delta_v(t)/\Delta_v(t^n)$ must be kept until local transfer
and the global product are both defined.

At $v\in S$ Part III, Theorem 9.1, gives the exact local
normalized orbital-integral identity
\begin{equation}\label{eq:localJO}
 J_{\tH_v}\bigl(z\,t^*,\mathcal O_{D,v}f_v\bigr)
 =b_{D,v}\,\widetilde\theta_v(z)^{-1}
 J_{H_v}(t,f_v),
\end{equation}
using the same quotient measures on the two copies of $H_v/Z_v$.
At $v\notin S$, $H_v\cong\GL_r(K_v)$ and the appropriate
root-sum comparison is not \eqref{eq:localJO}; it must use
Flicker--Kazhdan--Zou, and it can involve the kernel
$\mu_n^{k_v}/\mu_n$ of Proposition \ref{prop:power-etale}.
In particular the local factors $b_{D,v}$ for $v\in S$
\emph{cannot} simply be cancelled using
$\prod_v|n|_v=1$ without the split-side contributions.

\subsection{Haar ledger}\label{subsec:ledger}
For $v\in S$, Part II, Definition 2.1, fixes
\begin{equation}\label{eq:formaldegree}
 \vol(D_v^\times/F_v^\times;dg_{D,v})
 =d(\mathrm{St};dg_{\GL,v})^{-1},
\end{equation}
where $d(\mathrm{St};dg_{\GL,v})$ is the formal degree
relative to the chosen $\GL_r(F_v)/F_v^\times$ measure.
The probability measure used in Part III is
$d(\mathrm{St};dg_{\GL,v})\,dg_{D,v}$.
When \emph{both} sides of \eqref{eq:localJO} are scaled by the
same measure, $b_{D,v}$ is unchanged.  Different measures on the
$\GL$ side must still be reconciled with \eqref{eq:formaldegree}.

\begin{center}\small
\begin{longtable}{@{}p{0.27\linewidth}p{0.53\linewidth}p{0.10\linewidth}@{}}
\toprule
Item & Required statement or calculation & Status\\\midrule
\endhead
Ambient cocycle and common center & Part I, Lemma 5.1 and Prop. 5.5; restriction \eqref{eq:before} & [E]\\
Regular rational good locus & Theorem \ref{thm:globalgood} & [P]\\
Global rational/stable classes & Hilbert 90 and Kottwitz's framework; local embeddability into $D$ & [E]\\
Global torus coefficient & Same $dt$ in \eqref{eq:globalvol}, including $|W_\gamma|$ & [E]\\
Local $v\in S$ orbital transfer & Part III, Theorem 9.1, equation \eqref{eq:localJO} & [E]\\
Split-place $v\notin S$ invariance & Lemma \ref{lem:FKnecessary} (necessary); rank-two spherical calculation below & [P]\\
Split-place pointwise transfer and $\lambda_v$ & Proposition \ref{prop:conditionalglobal}; unproved existence & [C]\\
Unramified unit elements & \cite[\S\S11--12]{FK} under $(n,N_r)=1$; transport to Part I cover pending & [E]/[O]\\
Nonregular powers and singular germs & Remark \ref{rem:power-singular}; germ compatibility or support needed & [C]\\
Central identity class & Suppressed in Prop. \ref{prop:conditionalglobal}; separate germ/measure test below & [E]/[C]\\
Spectral expansion & Part I--II packet formula with restrictions; new factor matching & [E]/[C]\\
\bottomrule
\end{longtable}
\end{center}
Here $N_r$ is the least common multiple of the \emph{composite}
positive integers at most $r$ (the convention in [FK]), with $N_r=1$
when there are none.  The classical unramified unit matching in
\cite[\S12, Corollary]{FK} is established under $(n,N_r)=1$
and $p\nmid n$.  Thus $N_2=N_3=1$, and for prime $r=5$
we have $N_5=4$, automatically prime to the odd integer $n$.
Other ranks must retain the published hypothesis.

\subsection{Split-place invariance and a geometric comparison criterion}
At a place $v\notin S$ the algebra $D_v$ is split, but a rational
elliptic field $K[s]$ may decompose locally as a product of fields.
By Proposition~\ref{prop:power-etale}, its local $n$th-root fibre
modulo scalars can contain $n^{k_v-1}$ points.  The global class
$[s]$ supplied by Theorem~\ref{thm:globalgood} is nevertheless
unique.  Therefore a local sum over \emph{all} roots cannot be substituted
for the orbital integral of this one rational representative without
an additional compatibility statement.

\begin{definition}[FK-compatible ordinary functions]\label{def:FKcompatible}
Fix the local central characters $\theta_v,\widetilde\theta_v$
with $\theta_v(a)=\widetilde\theta_v(a^*)$.  A normalized
ordinary test function $f_v$ is called \emph{FK-compatible on the
regular-power locus} if, whenever $x,x'\in H_v$ have $x^n,x'^n$
regular and their star images satisfy
$({x'})^*=z x^*$ for a central lift $z\in\widetilde Z_v$
(after conjugating representatives), one has
\begin{equation}\label{eq:FKinvariance}
 J_{H_v}(x',f_v)
   =\widetilde\theta_v(z)^{-1}J_{H_v}(x,f_v).
\end{equation}
The case $x'=\xi x$ with \emph{scalar} $\xi\in\mu_n(F)$
is automatic, since $\theta_v(\xi)=1$.  The substantive
condition occurs on a product torus
$E_v=\prod_{i=1}^{k_v}E_i$ with $k_v\ge2$:
if $\zeta=(\zeta_i)\in\mu_n(E_v)\simeq\mu_n(F)^{k_v}$
is nonscalar and $x^n$ is regular, then
\begin{equation}\label{eq:nonscalar-root-invariance}
 J_{H_v}(\zeta x,f_v)=J_{H_v}(x,f_v).
\end{equation}
The odd-$n$ star is multiplicative on each lifted torus, and
$\zeta^*=1$, so $(\zeta x)^*=x^*$.  The kernel on the
quotient torus is $\mu_n(F)^{k_v}/\mu_n(F)$, of size
$n^{k_v-1}$ (Proposition~\ref{prop:power-etale}).
Equation \eqref{eq:FKinvariance} is the invariant-function
condition of \cite[\S6]{FK}, expressed with the present
inverse central transformation convention.
\end{definition}

\begin{lemma}[\Ptag\ Necessary split-place invariance]
\label{lem:FKnecessary}
Suppose a genuine $\widetilde f_v$ and a fixed nonzero constant
$\lambda_v$ satisfy
\begin{equation}\label{eq:splitpoint}
 J_{\widetilde H_v}(x^*,\widetilde f_v)
      =\lambda_v J_{H_v}(x,f_v)
\end{equation}
for all $x$ in a conjugation- and scalar-stable region where
$x^n$ is regular, with the identity at central translates understood
with the prescribed inverse-central-character law.  Then $f_v$ obeys
\eqref{eq:FKinvariance} on that region.
\end{lemma}
\begin{proof}
If $({x'})^*=z x^*$, conjugation invariance and the inverse
central transformation law of the genuine test function imply
$J_{\widetilde H_v}(({x'})^*,\widetilde f_v)
=\widetilde\theta_v(z)^{-1}
J_{\widetilde H_v}(x^*,\widetilde f_v)$.
Apply \eqref{eq:splitpoint} twice and divide by the fixed
nonzero $\lambda_v$.  In particular, for $x'=\zeta x$ with nonscalar
$\zeta\in\mu_n(E_v)$, one has $({x'})^*=x^*$ and obtains
\eqref{eq:nonscalar-root-invariance}.  The scalar-root case is
already automatic from the central-character convention.
Thus FK-compatibility is a \emph{necessary} condition,
not an asserted construction of $\widetilde f_v$.
\end{proof}

\begin{proposition}[\Ptag\ Rank-two tame spherical-unit check]
\label{prop:ranktwo-unit}
Let $F$ be nonarchimedean with $\mu_n\subset F$, $p\nmid n$,
$r=2$, and $K_0=\GL_2(\mathcal O_F)$.  Take a spherical ordinary
unit $f^0$, supported on $F^\times K_0$, with an unramified
central character, and the standard quotient measures. Then
$f^0$ obeys \eqref{eq:nonscalar-root-invariance} for all regular
split-torus $x$ with $x^n$ regular.  For a quadratic field torus
$E^\times$, the root-fibre condition modulo scalars is vacuous.
Moreover, when $f^0$ and the genuine spherical unit
$\widetilde f^0$ are transported from the standard
Flicker--Kazhdan cocycle and measure conventions,
\cite[\S12, Corollary]{FK} establishes their regular-power
matching with $\lambda_v=1$ in that normalization.  Whether the
same equality holds \emph{literally} for the selected ambient
cover and global measures of Parts I--IV requires their explicit
splitting and Haar identifications.
\end{proposition}
\begin{proof}
A split regular element is conjugate to
$\operatorname{diag}(a,b)$; multiplying by a scalar, first
consider $a,b\in\mathcal O_F^\times$.  The Iwasawa and
unipotent-orbit calculation, with $\operatorname{vol}(\mathcal
O_F)=1$, yields
\[
 O_{\GL_2}\!\left(\operatorname{diag}(a,b),1_{K_0}\right)
       =|a-b|_F^{-1},\qquad
 \Delta\!\left(\operatorname{diag}(a,b)\right)=|a-b|_F,
\]
so $J=1$.  Multiplication of either eigenvalue by a root of unity
preserves its valuation, hence the normalized orbital integral.
If the two eigenvalue valuations differ, neither $x$ nor such a
translate is conjugate into $F^\times K_0$; both orbital
integrals vanish. Scalar translations contribute only the
specified unramified central-character factor.  This proves the
ordinary split-torus invariance. For a quadratic field $E/F$,
Proposition~\ref{prop:power-etale} gives
$\ker(P_n:E^\times/F^\times\to E^\times/F^\times)=1$.
The final statement invokes the \emph{published} unit matching,
not the elementary split-torus calculation.  Its passage to the
specified ambient cover is a distinct normalization obligation.
\end{proof}

\begin{remark}[Scope of the published fundamental lemma]
\label{rem:fk-range}
The Flicker--Kazhdan result is a theorem for the matching
unramified Hecke units, not an existence result for arbitrarily
prescribed functions at finitely many exceptional places.
Its original hypotheses include $p\nmid n$ and
$(n,N_r)=1$. In rank two the only maximal-torus types are
split and quadratic field tori, so the ordinary invariance
calculation above is exhaustive.  At higher ranks, mixed tori
do not automatically invalidate spherical matching: the
Flicker--Kazhdan theorem covers them as well whenever its
hypotheses apply.
\end{remark}

\paragraph{A regularity and germ obligation.}
Equation~\eqref{eq:splitpoint} is imposed only where $x^n$ is
regular.  On a split torus, $x$ can be regular while $x^n$ is
singular (Remark~\ref{rem:power-singular}); the power map is a
local analytic isomorphism on the torus, but normalized orbital
integrals of a function at images approaching a singular element
may have nonconstant Shalika-germ behavior.  A full transfer
construction must either arrange support away from these
root-collision loci or prove the matching germ conditions there.
The same issue occurs at the complex places, although the cover
there splits: the ordinary power map can still carry regular
elements to singular ones. No universal vanishing of the germ
coefficients is claimed here.

\begin{proposition}[\Ptag\ Conditional geometric comparison with
an explicit normalization factor]\label{prop:conditionalglobal}
Fix the rational splittings and compatible central characters
of Parts I--III, and fix \emph{corresponding} local Haar measures,
global torus measures, and normalizer conventions on the ordinary
and covered sides.  Choose factorizable ordinary and genuine test
functions in a trace-class, regular elliptic support range for
which the two simple trace formulas apply. Assume:
\begin{enumerate}[label=(\roman*)]
\item For each $v\notin S$, $f_v$ is FK-compatible as in
Definition~\ref{def:FKcompatible}, and a genuine test function
$\widetilde f_v$ satisfies \eqref{eq:splitpoint} for every regular
rational local component $s_v$ of a globally good class $[s^n]$,
with a nonzero constant $\lambda_v$ independent of $s_v$ and of
its local torus type.  The prescribed central-translate identity
holds, $\lambda_v=1$ for almost all $v$, and the hypotheses include
the complex places.
\item For each $v\in S$, use the Part III, Theorem~9.1,
orbital-integral transfer $\widetilde f_v=\mathcal O_{D,v}f_v$,
with coefficient $b_{D,v}=|n|_v^{-(r-1)/2}$ and matching quotient
measures on $H_v/Z_v$ and $\widetilde H_v/\widetilde Z_v$.
The functions obey the reciprocal central transformation laws.
\item The common rational splitting and all local lifts satisfy
the star compatibility $\iota(s^n)=\iota(s)^n$; the global torus
coefficients (including $|W_s|$), central phases, and quotient
measures agree under the stipulated identifications.  The support
condition removes the central identity contribution.  The
ordinary and covered simple geometric expansions are absolutely
convergent in the selected test-function range.
\end{enumerate}
Put
\begin{equation}\label{eq:lambda-product}
 \boxed{\quad
 \Lambda=\Bigl(\prod_{v\notin S}\lambda_v\Bigr)
          \Bigl(\prod_{v\in S}b_{D,v}\Bigr).\quad}
\end{equation}
Then the regular elliptic geometric distributions, and hence the
corresponding simple spectral trace distributions, satisfy
\begin{equation}\label{eq:globalLambda}
 \boxed{\quad I_{\widetilde H}(\widetilde f)
              =\Lambda\,I_H(f).\quad}
\end{equation}
In particular the distributions agree if $\Lambda=1$.
\end{proposition}
\begin{proof}
Since the rational splitting $\iota$ is a homomorphism,
$\iota(s^n)=\iota(s)^n$; at every completion, $\iota(s)_v$
is a lift of $s_v$, so this equals its star lift $s_v^*$.
By Theorem~\ref{thm:globalgood}, globally good noncentral
rational classes are parametrized by the unique global root
class $[s]$ of $[s^n]$, preserving normalizer orders.
Every ordinary regular rational class $[s]$ is paired with its good
covered star class $[s^n]$; there is no additional hypothesis
that the ordinary class $[s]$ itself be good.  The cover has zero
regular orbital integrals at nongood lifts, which contribute
nothing to the covered geometric expansion.  Local compatibility
handles the multiple \emph{local} roots without introducing
additional rational summands.  Hypothesis (i) supplies
$\lambda_v J_{H_v}(s_v,f_v)$ off $S$, and (ii) supplies
$b_{D,v}J_{H_v}(s_v,f_v)$ on $S$. By (iii) the remaining torus
coefficients and central phases agree, and the product of Weyl
discriminants is $1$ by \eqref{eq:discprod}.  Every regular
rational term on the covered side is thus $\Lambda$ times the
corresponding ordinary term.  The central identity has been
suppressed, and the stipulated trace-class and convergence
conditions permit summation and the two trace formulas. This
gives \eqref{eq:globalLambda}.  The proof supplies \emph{no}
existence theorem for the matching functions or their exceptional
local constants.  It only establishes the implication under the
explicit hypotheses.
\end{proof}

\begin{remark}[Normalization versus canonical matching]
The condition $\Lambda=1$ is a normalization once nonzero
matching functions are known: rescaling one $\widetilde f_v$
rescales its $\lambda_v$ and hence $\Lambda$. It is \emph{not}
a theorem about canonically normalized units; rescaling a unit
changes its canonical value.  The meaningful unresolved question
is to compute the constants for the prescribed local transfers,
especially the finitely many exceptional split places, and to
verify the global product in those fixed conventions.
Likewise, compactness of $D^\times(K)Z(\mathbb A_K)\backslash
D^\times(\mathbb A_K)$ eliminates continuous spectrum but does
not by itself justify every unrestricted trace sum; the support
and trace-class conditions above are retained for this reason.
\end{remark}

\subsection{The identity class and an explicit normalization test}
\label{subsec:identity-test}
Because $r$ is prime and $D$ is division, every noncentral
rational element generates a degree-$r$ field and is regular
elliptic.  Thus, after quotienting by the rational center, the
central identity class is the only remaining singular rational
class.  Proposition~\ref{prop:conditionalglobal} suppresses it,
but it can be used as an additional test \emph{if} the test range
is enlarged and the necessary local germ identities are supplied.
\begin{remark}[Central discriminants versus orbital-integral germs]
\label{rem:identity-germs}
Let $d=r(r-1)/2$.  For any torus, as a regular $x\to1$
with $x^n$ regular, the \emph{discriminant ratio} satisfies
\[
 \frac{\Delta(x)}{\Delta(x^n)}\longrightarrow |n|_v^{-d},
\]
because $(1-u^n)/(1-u)\to n$ for each root value $u\to1$.
However this identity \emph{does not by itself yield}
$\widetilde f_v(1)=\lambda_v|n|_v^{-d}f_v(1)$
from \eqref{eq:splitpoint}.  On the split group, normalized
orbital integrals generally have nontrivial Shalika germs near
singular central elements; their limiting coefficients need not
be the evaluation of the test function at $1$.  The requisite
relation between \emph{central test-function values} is therefore
a separate local matching obligation, even when the regular
orbital integrals match.  Compare the unit-specific calculation
of Proposition~\ref{prop:ranktwo-unit} with the general
germ warning following Lemma~\ref{lem:FKnecessary}.
\end{remark}

At $v\in S$, Part III, Corollary~9.4 gives
$\widetilde f_v(1)=c_{D,v}f_v(1)$ with
$c_{D,v}=b_{D,v}|n|_v^{-r(r-1)/2}$.  If, as an \emph{additional unproved local identity-value
hypothesis}, every split place (including complex places) satisfies
$\widetilde f_v(1)=\lambda_v|n|_v^{-d}f_v(1)$ with
$d=r(r-1)/2$, and all denominators are nonzero, the ratio of
global identity-function values becomes
\begin{equation}\label{eq:identityproduct}
 \prod_v\frac{\widetilde f_v(1)}{f_v(1)}
   =\Lambda\prod_v|n|_v^{-r(r-1)/2}
   =\Lambda
\end{equation}
by the product formula.  This does not imply equality of the
\emph{identity trace-formula terms} without a separate choice of
matching quotient measures: homeomorphic adelic quotients need
not have equal volumes with arbitrarily scaled Haar measures.
Nor does it justify retaining the identity term in the simple
trace formula without checking its other hypotheses.

\begin{example}[A possible global numerical ledger]
\label{ex:ranktwoledger}
Let $K=\mathbb Q(\zeta_3)$, $r=2$, $n=3$, and choose a
quaternion division $K$-algebra ramified at the place over $3$
and at one other finite place; for instance the running example
$D=(-1,4-\zeta_3)_K$, ramified at $\mathfrak p_3$ and $\mathfrak p_7$.  At the wild division place,
$b_{D,v}=3$, $c_{D,v}=27$.  If the finite split places have
matching $\lambda_v=1$, the equation $\Lambda=1$ would require
$\lambda_\infty=1/3$ at the unique complex place.
Since the normalized complex absolute value has
$|3|_\infty=9$, the hypothesized central-value formula would then
give the complex identity factor $(1/3)/9=1/27$, balancing
$27$.  This is a \emph{necessary numerical consistency check}
under the stated assumptions, not a proof of matching complex
test functions or a global comparison.
\end{example}

\subsection{Why Li, Kottwitz, and DKV do not close the gap}\label{subsec:litroute}
The linear Jacquet--Langlands theorem of
Deligne--Kazhdan--Vign\'eras \cite{DKV} controls the ordinary
$D^\times$--$\GL_r$ relation, including the elliptic character
sign, but it does not automatically manufacture a genuine
representation of $\GDb_v$.  Kottwitz's analysis of stable
conjugacy and relative trace coefficients \cite{Kottwitz84}
provides the correct arithmetic setting; for the induced tori
$\operatorname{Res}_{E/K}\mathbb G_m$ one has
$H^1(K,T)=1$ by Shapiro's lemma and Hilbert 90, and embeddings
$E\hookrightarrow D$ are controlled by Brauer splitting.
Wen--Wei Li proves an invariant trace formula for adelic covers
under an invariant trace Paley--Wiener hypothesis for Levi subgroups
at \emph{real} archimedean places \cite{Li14}.  Here $K$ is totally
imaginary, so that real-place hypothesis is vacuous.  At the complex
places the specified covers split (Part I, Appendix~B, item~(R4);
see also its ambient-cover construction), and the linear invariant Paley--Wiener theorem
of Clozel--Delorme \cite{ClozelDelorme} supplies the relevant
archimedean background.  Li also discusses known real metaplectic
$\GL$ cases, but his general result does not assert the particular
\emph{equality} of ordinary-$D^\times$ and covered-$D^\times$
traces required in \eqref{eq:twocomparisons}.  The simpler already-proved trace formulas
of Parts I--II remain the appropriate initial comparison.

\section{Spectral separation and a research program}\label{sec:spectral}
\subsection{The actual spectral distribution identity}
Let $H=D^\times$ and $\widetilde H=\GDb$ globally, with
compatible central characters $\theta,\widetilde\theta$.
The quotient $H(K)Z_H(\A_K)\backslash H(\A_K)$ is compact
because $D$ is a division algebra over $K$; its covered analogue
is compact modulo the specified central preimage as well.
Thus the fixed-central-character right regular representations
have discrete Hilbert decompositions.  For factorizable test
functions with the inverse central transformation law, define
\begin{align}
 I_H(f)&=\sum_{\rho\in\mathcal A(H,\theta)}
       m_H(\rho)\prod_v\operatorname{tr}\rho_v(f_v),\\
 I_{\widetilde H}(\widetilde f)
     &=\sum_{\widetilde\rho\in\mathcal A(\widetilde H,\widetilde\theta)}
       m_{\widetilde H}(\widetilde\rho)
       \prod_v\operatorname{tr}\widetilde\rho_v(\widetilde f_v).
       \label{eq:spectralglobal}
\end{align}
Here $\mathcal A$ denotes the relevant automorphic spectrum
and $m$ is its actual spectral multiplicity, not an assumed $1$.
The expressions are interpreted on trace-class test functions
satisfying the simple trace-formula hypotheses; a bare
unrestricted infinite sum is not justified by notation alone.
Under the assumptions of Proposition~\ref{prop:conditionalglobal},
the comparison of geometric expansions supplies
$I_{\widetilde H}(\widetilde f)=\Lambda I_H(f)$;
ordinary equality holds when the fixed normalizations give
$\Lambda=1$.  It does not
identify summands.  This is a conditional identity of
\emph{distributions}, with both the multiplicities and local
trace factors retained.

\subsection{Individual local characters and integrality}
An equality of global spectral sums is a statement about
\emph{distributions}.  To isolate a single local irreducible at a
wild place, one needs an argument that the relevant coefficients
are integers (or multiplicities which can be shown to isolate a
single constituent).  Linear independence alone can separate
finite lists of characters; it does not supply the required
integrality from a unitary Fourier transform.

\begin{remark}[What an integer spectral coefficient would prove]
If a justified global-to-local argument supplies
$a_{\rho\tau}\in\mathbb Z$ for every $\tau$ in the row of a fixed
$\rho$, Proposition~\ref{prop:int} (Part III, Corollary~5.2)
forces $L_{D,v}\Theta_\rho$ to be one irreducible genuine
character with positive sign.  If integrality holds for all
rows, unitarity implies a bijection.  This is an application of
an already proved lattice criterion, not a second conditional
theorem; the missing step is precisely the integrality
hypothesis at a wild place.
\end{remark}

\paragraph{Separating the remaining tasks.}
First, Theorem~\ref{thm:quad-toral} proves toral permutation
positivity in rank two for every odd covering degree.  The separate
rank-two project \cite{RankTwo0} addresses the remaining nonabelian
gluing problem and must justify the choice of a finite level before
any finite character calculation becomes decisive.  Second, at split
places the root sum in Zou's (3.2) should be written against the
kernel and index in Proposition \ref{prop:power-etale}, including
the chosen Haar measures and central phases.  Third, one needs the explicit ambient-cover identification
that transports the known unramified unit matching of
\cite[\S12]{FK} to the selected convention, together with
functions at the finitely many exceptional split places and a
complete global coefficient identity as in
Proposition~\ref{prop:conditionalglobal}. Every ordinary rational regular class is paired with the star of
its global representative; the ordinary class itself need not be
good.
Finally, one must justify globalization and spectral separation
for representations of the covered inner form, and remove the
auxiliary supercuspidal and non-cancellation hypotheses of Parts
I--II.  These are mathematical work packages, not restatements of
an already-proved theorem.

\section{Selected glossary}\label{sec:glossary}
\begin{description}[style=nextline,leftmargin=3.0cm]
\item[Good element] A regular lift commuting with the full preimage of its ordinary centralizer; genuine invariant characters vanish on the nongood regular locus.
\item[Scalar-central cover] A cover whose entire inverse image of $F^\times I_r$ is central; here forced by $n\mid r-1+2rc_v$.
\item[Power-map kernel] $\ker(P_n)$ on $T(F)/F^\times$, distinct from the image index.  The kernel can be nontrivial on product tori even when the corresponding field-torus map is injective.
\item[Star map] $h^*=\widetilde h^{\,n}$, independent of lift; multiplication on tori uses $n$ odd.
\item[Weyl discriminant] The factor $\Delta=|\det(1-\operatorname{Ad}(t)|\mathfrak g/\mathfrak t)|^{1/2}$; it is not the unnormalized characteristic-polynomial discriminant.
\item[Toral Fourier test] Expand the discriminant ratio $R_E$ in characters of the compact abelian group $E^\times/F^\times$.  Integral positive coefficients ensure torus-restriction positivity, not full-group integrality.
\item[Nonabelian gluing] The additional requirement that compatible genuine covered-torus characters arise from an integral character of the full group $\widetilde{D^\times}$.
\item[Twisted finite quotient] A finite quotient of $H/Z$ with the central character retained as a multiplier, used to compute projective/genuine irreducible characters.
\item[Packet identity] An identity of traces over collections of automorphic representations; it does not assert a bijection between individual local irreducibles.
\item[Formal degree] The Plancherel normalization of a square-integrable representation relative to a specified quotient Haar measure.
\item[Root-sum relation] Zou's equation (3.2), which sums over all $n$th roots satisfying the stated regularity condition rather than choosing one arbitrarily.
\end{description}

\section*{Acknowledgments}
This Part~IV manuscript was prepared with mathematical and editorial
assistance from ChatGPT (OpenAI).  Claude (Anthropic) supplied
referee reports on Drafts~1--3, including the correction to the
Haar-modulus proof and the split-place matching obligations.  Earlier Parts I--III were developed
with separate assistance acknowledged there.  The author contributions and the referee corrections are documented
in the separate revision log.  The quadratic toral theorem is included with proof in this
concluding draft; the detailed wild rank-two calculations now form
a standalone Draft~0 companion project.  The remaining research problems are
recorded in the status table and Section~\ref{sec:spectral}.

\begin{thebibliography}{99}
\bibitem{DKV} P. Deligne, D. Kazhdan and M.-F. Vign\'eras,
\emph{Repr\'esentations des alg\`ebres centrales simples $p$-adiques},
in \emph{Repr\'esentations des groupes r\'eductifs sur un corps local},
Travaux en Cours 8, Hermann, Paris (1984), 33--117.
\bibitem{FK} Y. Flicker and D. Kazhdan,
\emph{Metaplectic correspondence}, Publ. Math. IH\'ES \textbf{64}
(1986), 53--110, especially \S\S2--4, 12, 18, 24--28.
\url{https://www.numdam.org/articles/10.1007/BF02699192/}.
\bibitem{Zou24} J. Zou,
\emph{Local metaplectic correspondence and applications},
arXiv:2209.05852v3 (19 December 2024), especially
(2.2), (3.2), Remarks 3.1 and 3.8, and Proposition 3.6.
\url{https://arxiv.org/abs/2209.05852}.
\bibitem{KP} D. Kazhdan and S. Patterson,
\emph{Metaplectic forms}, Publ. Math. IH\'ES \textbf{59}
(1984), 35--142.
\bibitem{Kottwitz84} R. E. Kottwitz,
\emph{Stable trace formula: cuspidal tempered terms}, Duke Math. J.
\textbf{51} (1984), 611--650.
\bibitem{ClozelDelorme} L. Clozel and P. Delorme,
\emph{Le th\'eor\`eme de Paley--Wiener invariant pour les groupes de Lie
r\'eductifs}, Invent. Math. \textbf{77} (1984), 427--453.
\bibitem{Li14} W.-W. Li,
\emph{La formule des traces pour les rev\^etements de groupes
r\'eductifs connexes. IV. Distributions invariantes},
Ann. Inst. Fourier \textbf{64} (2014), 2379--2448.
\url{https://doi.org/10.5802/aif.2915}.
\bibitem{JC} D. Joyner with Claude (Anthropic),
\emph{On the metaplectic analog of Kazhdan's ``endoscopic'' lifting},
corrected 2026 version of the 1988 paper,
\texttt{metaplectic\_kazhdan\_revised12.tex}, especially \S\S2, 6.
\bibitem{JoynerI} D. Joyner, Claude (Anthropic), and ChatGPT (OpenAI),
\emph{A metaplectic Jacquet--Langlands transfer for division algebras
of prime degree, I: Local matching and a global transfer under
elliptic-support conditions}, Draft 8 (2026), especially \S\S2--6.
\bibitem{JoynerII} D. Joyner, Claude (Anthropic), and ChatGPT (OpenAI),
\emph{A metaplectic Jacquet--Langlands transfer for division algebras
of prime degree, II: Matching of arbitrary functions and an improved
global transfer}, Draft 8 (2026), especially \S\S2, 5--7.
\bibitem{RankTwo0} D. Joyner, ChatGPT (OpenAI), and Claude (Anthropic),
\emph{Wild-place metaplectic character transfer in degree two:
Toral positivity and the nonabelian gluing problem},
Draft 0 (2026), companion manuscript in the Part-IV folder.
\bibitem{JoynerIII} D. Joyner, ChatGPT (OpenAI), and Claude (Anthropic),
\emph{A metaplectic Jacquet--Langlands transfer for division algebras
of prime degree, III: The division-algebra side and its local
character transform}, Draft 9 (2026), especially Corollary~7.8
and \S\S3--5, 7--9, 11--12.
\end{thebibliography}
\end{document}
